The function $\varphi(y) = y/(1 + y^{2})$ satisfies $|\varphi(y)| \leq \tfrac12$, since $1 + y^{2} - 2|y| = (1 - |y|)^{2} \geq 0$, with equality at $y = \pm 1$; and $\varphi'(y) = (1 - y^{2})/(1 + y^{2})^{2}$, so $\varphi'(0) = 1$. Hence $\sup_{x}|a_{\tau}(x)| = \tfrac12\tau^{1/2}$, attained at $x = \tau^{2}$, which tends to zero uniformly in $x$; and $\sup_{x}|a_{\tau}|/\tau = \tfrac12\tau^{-1/2}$, whose integral over $(0, \tau_{0}]$ is $\tau_{0}^{1/2}$ (substitute $\tau = s^{2}$). With $\|B_{2}\| = 1$ these are the decay clauses of Definition \claimref{obj.interior.weighted-blow-up-hubble-normalisation-admissible} for $E_{\tau} = a_{\tau}B_{2}\otimes\hat\theta^{2}$, uniformly on the patch. By the chain rule $\partial_{x}a_{\tau}(x) = \tau^{-3/2}\varphi'(x/\tau^{2})$, which is $\tau^{-3/2}$ at $x = 0$; the check confirms it at $\tau = n^{-2}$ through the exact difference quotient $a_{\tau}(h)/h = n^{3}/(1 + (h/\tau^{2})^{2})$. If $x$ is the coordinate of axis $j$, the unit vector is $\tau^{-p_{j}}\partial_{x}$ and a Hubble-normalised length is $\tau^{-1}$ times a proper length, so the derivative along unit tangential vectors per unit Hubble-normalised length is $\tau^{1 - p_{j}}\partial_{x}$, giving $\tau^{-1/2 - p_{j}}$ at $x = 0$; on the constraint variety $p_{j} \geq -\tfrac13$, so the exponent is at most $-\tfrac16$ and the derivative diverges for every axis. This is the normalisation of the conditional bound filed alongside, synthesis.tangential-derivative-transport-bound.

Along the curve in the direction $\hat\theta^{2}$ through a point with coordinate $x$ (on which $x$ is constant), the Hubble-normalised generator is $G = (p_{2} + a_{\tau}(x))B_{2}$, constant along the curve, and the model generator is $H = p_{2}B_{2}$. Both are multiples of $B_{2}$, so the transports are $T(x) = \exp(\lambda(p_{2} + a_{\tau}(x))B_{2})$ and $T_{0} = \exp(\lambda p_{2}B_{2})$, and $\|T - T_{0}\| \to 0$, as Theorem \claimref{interior.transport-level-convergence-along-tangential-curves} guarantees from the decay clauses alone. Since $B_{2}^{2} = 1$ (the boost $e_{0}e_{2}$ of $\Cl(1,3)$ squares to one), $\exp(\sigma B_{2}) = \cosh\sigma + B_{2}\sinh\sigma$, and $\partial_{x}T = \lambda\,\partial_{x}a_{\tau}(x)\,B_{2}T$ because everything commutes. Now $B_{2}\exp(\sigma B_{2}) = \sinh\sigma + B_{2}\cosh\sigma$ has $B_{2}$-coefficient $\cosh\sigma \geq 1$, so in any norm dominating the coefficient of $B_{2}$ (the Clifford coefficient norm, or the operator norm on the six-dimensional tractor representation, where $B_{2}$ acts with eigenvalue one) $\|\partial_{x}T(0)\| \geq \lambda\,\tau^{-3/2}$, and $\geq \lambda\,\tau^{-1/2 - p_{j}}$ in the Hubble-normalised derivative. The check verifies, in signature $(-,+,+,+)$ (the square is the same in $(+,-,-,-)$), $B_{2}^{2} = 1$ and the identity $B_{2}(c + sB_{2}) = s + cB_{2}$ on rational points $c^{2} - s^{2} = 1$ of the hyperbola. Since the example satisfies the decay clauses and the Kasner constraint, no consequence of those hypotheses can bound $\partial_{x}E_{\tau}$ or $\partial_{x}T$ uniformly as $\tau \to 0$. The Einstein clause is the only remaining hypothesis of admissibility; this family is an abstract error field, not claimed to arise from any metric, and whether the Einstein clause supplies a bound is the open question.
